\chapter*{Abstract}
\label{summary}

Adolescent mental health support remains fragmented: self-care tools including mood or sleep trackers such as Daylio, Wysa and Bearable, rarely connect to licensed professional care, while therapy platforms such as BetterHelp and Talkspace typically specialize in only finding you a trusted therapist without giving the therapist enough context as to how you are doing in order to efficiently conduct meaningful sessions. And while there are platforms for people to have deep conversations with close friends and family, it is hard for parents to actually know how their children are doing and feeling. And as a trend of young people seeking mental health advice from AI, a centralized platform where the user can choose to share their daily reflective tracking data with AI would be more beneficial in case they truly want a deeply contextual advice instead of a vague one. This gap motivates a platform that unifies these fragmented features so that the user don't have to spend unnecessary time explaining themselves every time they want to use  a feature from a different platform, providing the best mental health support as we can.

This thesis presents uMatter, a mental health support platform for teenagers that combines self-care features, an AI companion, peer social features as well as professional therapy features.

uMatter is implemented as a microservices system: seven Spring Boot services (authentication, tracking, AI, dashboard, therapist, social, and notification) sit behind an Nginx API gateway, communicate asynchronously through RabbitMQ, and follow a database-per-service pattern backed by PostgreSQL, Redis, and MinIO. Access control uses stateless JWT authentication with a JWKS-published RSA key and a permission-based data-sharing model so users can selectively grant professionals access to their tracked data. The AI companion is grounded in each user's own tracking history and served through Google's Gemini 2.5 Flash. The system was hardened through structured security audits covering injection risks, insecure direct object references, and mass assignment, and a database-level atomic slot claim, backed by a unique constraint, was introduced to close a race condition in appointment booking. The full stack was deployed and validated on cloud infrastructure, running as 20 containers.

The resulting platform demonstrates that a unified self-care and professional therapy experience can be successfully delivered through a single, security-conscious microservices architecture. However, as an academic prototype, uMatter scopes out key real-world operational challenges. It excludes legal and compliance frameworks, dispute resolution mechanisms for handling user or therapist conflicts, and the monetization model essential for incentivizing licensed professionals. Despite these out-of-scope administrative limitations, this thesis establishes a robust technical foundation. Future work can build upon this architecture by integrating monetization systems, comprehensive legal safeguards, clinical collaboration requirements, and layered risk detection, ultimately advancing uMatter toward a fully viable ecosystem for adolescent mental health support.